\chapter{Introduction}
\label{chap:introduction}

Les objets connectés et les systèmes industriels forment des réseaux
hétérogènes dans lesquels une passerelle, un courtier MQTT, un automate ou un
service web peut devenir une étape d'un chemin d'attaque. Évaluer leur sécurité
suppose de découvrir la topologie, de choisir des actions adaptées aux
protocoles, puis de vérifier les hypothèses par des observations et d'en
conserver les preuves dans des artefacts comparables entre plusieurs campagnes.
Les graphes d'attaque offrent un cadre pour représenter les dépendances entre
connectivité, configuration et compromission\cite{ou2005mulval}.

Le projet LANCE (\emph{LLM Agent for Network Compromise Evaluation}) étudie
l'emploi de \glspl{llm} pour automatiser cette démarche dans des réseaux
\gls{iot}/\gls{ot}. Son agent conduit une campagne structurée, depuis la
découverte du réseau jusqu'au rapport. Cette application pose une difficulté
centrale~: un modèle généraliste peut produire une réponse plausible sans
respecter le rôle attendu, le format d'un appel d'outil ou le niveau de preuve
nécessaire à une conclusion de sécurité.

La mission principale du stage a consisté à concevoir une architecture agentique
de type \gls{hmoe}, fondée sur un modèle de base partagé et sur des experts
spécialisés par \gls{qlora}. Chaque expert apprend un rôle opérationnel
organisation de la campagne, reconnaissance, analyse de vulnérabilités ou
exploitation et un routeur sélectionne l'adaptateur correspondant à la phase
exécutée. Le travail couvre ainsi la préparation de données issues de traces,
leur contrôle, l'entraînement, le service des adaptateurs et leur
intégration au cycle d'action de l'agent.

Le service HMoE est déployé localement derrière une API compatible OpenAI. Le
pipeline peut ainsi l'utiliser comme un fournisseur de modèle sans exposer la
sélection interne des experts.

Le pipeline d'exécution et le harness expérimental constituent le support de
cette spécialisation. Le premier impose les transitions, les outils autorisés
et les contrats de sortie ; le second déploie des scénarios reproductibles,
conserve la vérité terrain du côté de l'évaluateur et mesure le comportement
observé pendant une campagne. Leur consolidation a porté sur la robustesse de
l'orchestration, la collecte de preuves, la gestion des exécutions et
l'évaluation déterministe et sémantique.

La problématique peut dès lors être formulée ainsi~:
\emph{comment spécialiser un agent fondé sur un LLM local pour les différentes
étapes d'un audit IoT/OT, puis démontrer la qualité de son comportement au moyen
d'un protocole d'exécution et d'évaluation reproductible ?}

Le \cref{chap:contexte} présente le cadre et les objectifs de la mission. Le
\cref{chap:implementation} décrit l'intégration de la solution, notamment le
pipeline et le harness. Enfin, le \cref{chap:science} expose \gls{qlora}, le
routage \gls{hmoe} et leur évaluation.
